\documentclass[a4paper]{article}

\usepackage{graphicx}
\usepackage{amsmath,amssymb}
\usepackage{minted}
\usepackage{multirow}
\usepackage{float}
\usepackage{todonotes}
\usepackage{forest}
\usepackage{xstring}
\usepackage{xcolor}
\usepackage[most]{tcolorbox}
\usepackage{xparse}
\usepackage{natbib}

\usetikzlibrary{graphs,graphdrawing}
\usegdlibrary{layered}

% for external graphviz
\input{/home/donkarlo/Dropbox/repo/nd_ai_project/src/nd_ai/knoweldge_representation/language/formal/computer/programming/tex/latex/code_clips/external_graphviz.tex}

% for tags
\input{/home/donkarlo/Dropbox/repo/nd_ai_project/src/nd_ai/knoweldge_representation/language/formal/computer/programming/tex/latex/parts/control_sequence/kind/semtag/semtag.tex}

% for links
\input{/home/donkarlo/Dropbox/repo/nd_ai_project/src/nd_ai/knoweldge_representation/language/formal/computer/programming/tex/latex/code_clips/seehere.tex}

\title{Spline interpolation}
\date{}
\author{Mohammad Rahmani}

\begin{document}
\maketitle
    \seeherefooter{https://en.wikipedia.org/wiki/Spline_interpolation}
    
    \cite{lovegrove-2013-spline-fusion-a-continuous-time-representation-for-visual-inertial-fusion-with-application-to-rolling-shutter-cameras} propose a continuous-time spline-based framework for visual-inertial SLAM and calibration. Instead of representing the sensor trajectory only at discrete timestamps, they model it as a continuous curve using cubic B-splines. This allows measurements from unsynchronized and high-rate sensors, such as cameras and IMUs, to be fused in a common temporal framework. The method is also useful for rolling-shutter cameras, because each image scanline can be associated with its own measurement time.





    \bibliographystyle{plainnat}
    \bibliography{/home/donkarlo/Dropbox/repo/nd_shared_research/bibliography/bibliography.bib}
\end{document}